\section{Bugünkü durum}\label{sec:bugun}
\noindent{\sffamily\small\color{muted}27 Ekim 2026, bu raporun yazıldığı an}\par\smallskip

\subsection{Kart tasarımındaki modeller}

Tablo~\ref{tab:kart}, kart tasarımına giren \gBoardModels{} modelin bugünkü durumunu tek yerde toplar. “Saat” sütunu, açık kaynak akışın zamanlamayı tutturduğu çekirdek clock'tur; token/s bu saatten hesaplanır. Henüz yerleştirilmemiş modeller “sürüyor” olarak görünür; tablo \code{veri.py} yeniden çalıştırılınca güncellenir.

\begin{table}[htb]
\centering\small
\setlength{\tabcolsep}{4.5pt}
\caption{\gBoardModels{} model, Nexys A7-100T kart tasarımında. cycle/token kart simülasyonundan (otomatik ilk koşunun ortalaması); “Sim.” tam kart simülasyonunun Python referansıyla karşılaştırması.}\label{tab:kart}
\begin{tabular}{@{}lllrrlrl@{}}
\toprule
\thead{Model} & \thead{Veri} & \thead{Weight} & \thead{cycle/token} & \thead{token/s, 50 MHz} & \thead{Saat} & \thead{token/s, saatte} & \thead{Sim.} \\
\midrule
\tabinput{data/kart.tex}
\bottomrule
\end{tabular}
\end{table}

\begin{figure}[htb]
\centering
% GENERATED by report/gelisim/veri.py
\begin{tikzpicture}
\begin{axis}[width=\linewidth, height=6.2cm, ybar, bar width=7pt, ymode=log, log origin=infty,
  ymin=80, ymax=120000, ylabel={cycle / token}, xtick=data, xticklabel style={rotate=60, anchor=east, font=\sffamily\scriptsize},
  symbolic x coords={bigram,mlp,rnn,gru,bow,gpt2,microgpt,llama,nanochat,stories260k,tsgpt2,tsgru,tsllama,tsmicrogpt,tsnanochat,gpt2bin,mlplb25p1,stories260kbin},
  xticklabels={{bigram},{mlp},{rnn},{gru},{bow},{gpt2},{microgpt},{llama},{nanochat},{stories260k},{ts\_gpt2},{ts\_gru},{ts\_llama},{ts\_microgpt},{ts\_nanochat},{gpt2-bin},{mlp-lb25p1},{stories260k-bin}},
  ymajorgrids, enlarge x limits=0.03, axis x line*=bottom, axis y line*=left]
\addplot[fill=accent!75, draw=none] coordinates {(bigram,157) (mlp,6722) (rnn,6464) (gru,6605) (bow,6368) (gpt2,8712) (microgpt,8746) (llama,9490) (nanochat,10813) (stories260k,50397) (tsgpt2,41781) (tsgru,30407) (tsllama,50397) (tsmicrogpt,40570) (tsnanochat,50680) (gpt2bin,4983) (mlplb25p1,1874) (stories260kbin,27992)};
\end{axis}
\end{tikzpicture}

\caption{Token başına cycle (logaritmik eksen). names modelleri birkaç bin, TinyStories modelleri 30–50 bin cycle; düşük bitli modeller INT8 karşılıklarının altında.}
\label{fig:cyc}
\end{figure}

\subsection{Ne bitti, ne kaldı?}

\begin{table}[htb]
\centering\small
\caption{Projenin durumu.}\label{tab:durum}
\begin{tabularx}{\textwidth}{@{}>{\raggedright\arraybackslash}p{4.6cm}>{\raggedright\arraybackslash}p{2.2cm}L@{}}
\toprule
\thead{İş} & \thead{Durum} & \thead{Not} \\
\midrule
Integer aritmetik, ispatlı sınırlar & \textcolor{good}{bitti} & LN, RMSNorm, RoPE, QK-norm, GRU durumları \\
9 mimari + stories260K, tek engine & \textcolor{good}{bitti} & PYNQ dönemi engine'i, $P=8/16/32$ \\
Kart tasarımı (Nexys A7-100T) & \textcolor{good}{bitti} & sampling, detokenizer, VGA, kart üstü kontrol \\
Düşük bitli weight'ler & \textcolor{good}{bitti} & GEMVB + SCALE; 3 model kartta \\
Tam kart simülasyonu & \textcolor{good}{\gBoardPass{}/\gBoardModels{}} & token, hash, ekran \\
Rastgele / köşe / 4-state & \textcolor{good}{bitti} & \gFuzzRuns{} koşu, \gDirectedCases{} durum \\
Sentez sonrası netlist simülasyonu & \textcolor{warm}{kısmi} & gru, gpt2-bin, mlp-lb25p1 kart düzeyinde; 4 model engine düzeyinde \\
Yerleşim, yönlendirme, Fmax & \textcolor{warm}{\gPnrDone{}/\gPnrTotal{}} & büyük modeller 25 MHz'te \\
Fiziksel kart & — & kullanıcının kararıyla yapılmayacak \\
\bottomrule
\end{tabularx}
\end{table}

\subsection{Sonraki adımlar}
\begin{enumerate}
  \item \textbf{Engine'i pipeline'a bölmek, 50 MHz:} NORM çıkışı gibi tek cycle'lık eleman bazlı yollara bir register eklemek (§\ref{sec:e8}). Bedeli birkaç cycle; kazancı büyük modellerde yaklaşık iki kat token/s. Bütün testlerin yeniden geçmesi gerekir.
  \item \textbf{Yerleştirilen netlist'in simülasyonu, 18 model:} akışın yerleştirdiği netlist'i Verilog olarak kaydedip gerçek clock ayarı ve gerçek debounce süreleriyle tam kart simülasyonundan geçirmek (§\ref{sec:e9}).
  \item \textbf{minGPT'nin sıralama modeli, 19. model:} projenin başladığı \code{gpt-nano} sıralama görevi bugün kart tasarımında yok; o model eski PYNQ hızlandırıcısında çalışıyordu. Zoo'daki \code{gpt2} aynı GPT-2 bloğu ama names'te eğitildi. Ekrana girdi dizisini ve sıralanmış çıktıyı yazan bir görünümle eklenebilir. Sunumun konusu olduğu için önemli.
  \item \textbf{Daha uzun binary QAT (stories260K)} ve \textbf{iki yollu LittleBit'in donanımda toplanması}.
  \item \textbf{Yosys DSP hatasının} küçük örnekle (\code{requant\_netlist\_check.sh}) Yosys geliştiricilerine bildirilmesi; düzelirse DSP'li ve bit-exact bir açık kaynak bitstream.
  \item \textbf{Sunumun geri kalanı:} 16 Ekim'deki sunum yalnız minGPT kısmını kapsıyordu; implementasyonu anlatan yeni bir sunum, bu raporla birlikte hazırlanıyor.
\end{enumerate}

\subsection{Geriye bakış}

On beş günde projenin üç kez yön değiştirdiği görülüyor: tek modelden aileye, PYNQ'dan Nexys'e, bilgisayarın sürdüğü hızlandırıcıdan kartta kendi başına çalışan modele. Her yön değişikliği kullanıcıdan geldi ve her biri iddiayı daha dürüst hâle getirdi. Değişmeyen şey bit-exact kuralıydı; bu kural sayesinde her yeni katman (sampling, ekran, düşük bit, yeni araçlar) bir öncekinin doğruluğunu bozmadan eklenebildi. Üç genel ders:

\begin{itemize}
  \item \textbf{Spesifikasyon çalıştırılabilir olmalı.} Python'daki integer model hem bir referans hem de bir test kâhinidir; RTL'deki her hata ve araçlardaki üç hata ona karşı bulundu.
  \item \textbf{Araçlar da hata yapar.} Aynı tasarımı birbirinden bağımsız zincirlerden geçirmek, tek bir zincirin göremeyeceğini gösterir.
  \item \textbf{Ölçülmeyen sayı tasarımın zayıf yanını saklar.} Fmax ölçülünce tek cycle'lık yolların bedeli görüldü; sıradaki iş tam olarak bu.
\end{itemize}
